\section*{HISTORICAL NOTE (Chapters I through IV)}
\phantomsection
\addcontentsline{toc}{section}{Historical Note on Chapters I--IV}

``Aliquot selectos homines rem intra quinquennium absolvere posse puto'' Leibniz ({[}12 b{]}, vol.~7, p.~187).

(Numbers in brackets refer to the bibliography at the end of this Note.)

The study of what is usually called the ``foundations of mathematics'', which has been unflaggingly pursued since the beginning of the nineteenth century, could not have achieved success without a parallel effort of systematization of logic, or at least of those parts of logic which govern the relationships of mathematical propositions. Thus the history of the theory of sets and the formalization of mathematics cannot be separated from that of ``mathematical logic''. But traditional logic, like that of the modern philosophers, extends in principle over a field of applications far wider than mathematics. The reader will therefore not find a history of logic, even in a summary form, in this Note; we have limited ourselves, so far as possible, to retracing the evolution of logic only in so far as it has influenced the evolution of mathematics. Thus we shall say nothing about non-classical logics (many-valued logics, modal logics); a fortiori, we shall not go into the history of the controversies which, from the time of the Sophists to that of the Vienna school, have divided philosophers as to the possibility and manner of applying logic to the objects of the external world or the concepts of the human mind.

Today it is well established that mathematics was already highly developed in pre-Hellenic times. Not only are the (extremely abstract) notions of integers and measurement of magnitudes freely used in the most ancient documents which have come down to us from Egypt and Chaldaea, but the algebra of the Babylonians, by the elegance and sureness of its methods, cannot be regarded as simply a collection of problems solved by empirical devices. And, although there is nothing in the texts resembling a ``proof'' in the formal sense of the word, we may justly consider that the discovery of such methods of solution, whose generality is apparent from the particular numerical applications treated, could not have been achieved without a minimum of logical reasoning (perhaps not entirely conscious, but rather of the type on which a modern algebraist relies when undertaking a calculation, before writing out all the details) ({[}1{]}, pp.~203 ff.).

The originality of the Greeks consists precisely in a conscious effort of arranging mathematical proofs in a sequence so that the passage from one step to the next leaves nothing in doubt and compels universal assent. Of course, the Greek mathematicians, just like their present-day successors, made use of ``heuristic'' rather than rigorous arguments in the course of their researches; for example, Archimedes, in his Treatise on Method {[}4 bis{]}, refers to results ``found, but not proved'' by earlier mathematicians (*). But, from the earliest detailed texts known to us (which date from the middle of the fifth century B.C.) the ideal ``canon'' of a mathematical text is fixed. It found its complete realization in the works of the great masters Euclid, Archimedes, and Apollonius. Their notion of proof differs in no respect from ours.

We do not possess any text which would allow us to trace the first steps of this ``deductive method''. At its first known appearance, it is already nearly perfected. We can only think that it was a natural development from the perpetual search for ``explanations'' of the world, which characterized Greek thought from the time of the Ionian philosophers of the seventh century B.C. onward; furthermore, tradition is unanimous in ascribing the development and perfection of the method to the Pythagorean school, at a period somewhere between the end of the sixth century and the middle of the fifth.

This ``deductive'' mathematics, precise in its aims and rigorous in its methods, was to occupy the philosophical and mathematical reflections of the succeeding ages. We see, on the one hand, the edifice of ``formal'' logic slowly being built up, on the model of mathematics, and culminating in the creation of formalized languages; and on the other hand, principally from the beginning of the nineteenth century onward, especially after the advent of the theory of sets, the concepts which lie at the root of mathematics being more and more closely questioned and examined in an effort to clarify their nature.

